\section{Appendix}\label{sec:appendix}

% ─────────────────────────────────────────────────────────
\subsection{Detailed Simulation Flow Chart}\label{sec:app_flow}

Figure~\ref{fig:flow_detailed} expands the high-level process diagram
from the Methodology section to show every resource request, branching
decision, and stochastic sampling step in the model.
Starting from a kit order's arrival, the chart follows operator assignment,
per-item pick routing (Kardex picks, manual hand-reach picks, and
reach-truck-required picks), resource queuing for both the operator pool
and the RT fleet, the RT return-to-depot leg, and final lead-time recording.
Shift-break logic and the batch-arrival driver are also shown.

\begin{figure}[H]
  \centering
  \includegraphics[width=\textwidth]{sim_flowchart_detailed}
  \caption{Detailed simulation flow chart.  The diagram expands
    Figure~\ref{fig:simflow} in the Methodology with all resource-request nodes,
    stochastic-sampling symbols (order inter-batch gap, batch size,
    items per kit, lognormal pick durations), and the reach-truck
    return leg that keeps the RT resource busy until the truck is
    back at the depot.}
  \label{fig:flow_detailed}
\end{figure}

% ─────────────────────────────────────────────────────────
\subsection{Full KPI Table}\label{sec:app_kpi_full}

\begin{table}[htbp]
\centering
\caption{Full KPI table, mean over $N=20$ replications.}
\label{tab:kpi_full}
\scriptsize
\setlength{\tabcolsep}{3pt}
\begin{tabular}{l*{6}{c}}
\toprule
KPI & \makecell{Current\\SAP} & \makecell{Capacity\\heuristic} & \makecell{Usage\\ABC} & \makecell{Double\\ABC} & \makecell{Travel-dist.\\opt.} & \makecell{Line-aware\\slotting} \\
\midrule
Lead time (min) & 12.37 & 14.71 & 15.24 & 14.74 & 24.60 & 6.91 \\
Prep time (min) & 3.92 & 4.31 & 4.43 & 4.34 & 5.62 & 2.65 \\
Total wait (min) & 8.45 & 10.40 & 10.82 & 10.40 & 19.03 & 4.26 \\
P95 lead time (min) & 35.9 & 43.0 & 42.9 & 41.9 & 70.6 & 18.8 \\
Walk distance (m) & 89.0 & 106.1 & 108.8 & 104.6 & 96.1 & 92.4 \\
Throughput (orders/day) & 398.5 & 398.4 & 398.4 & 398.4 & 396.2 & 398.6 \\
RT utilization & 0.217 & 0.228 & 0.244 & 0.246 & 0.440 & 0.013 \\
Operator utilization & 0.405 & 0.446 & 0.457 & 0.448 & 0.576 & 0.275 \\
Kardex utilization & 0.079 & 0.079 & 0.079 & 0.078 & 0.078 & 0.079 \\
Orders completed (5 days) & 1943 & 1942 & 1942 & 1942 & 1931 & 1943 \\
\bottomrule
\end{tabular}
\end{table}


Table~\ref{tab:kpi_full} reports the mean of all ten tracked KPIs over
$N = 20$ paired replications for each policy.
Lead time, total waiting time, RT utilization, and throughput are the
primary comparison metrics discussed in the Results section.
The remaining rows --- prep time, P95 lead time, walking distance,
operator utilization, Kardex utilization, and orders completed over the
five-day run --- are included for completeness and to let readers
cross-check the results in Minitab using the per-replication export
(\texttt{kpi\_by\_replication.csv}).

% ─────────────────────────────────────────────────────────
\subsection{Model Assumption Register}\label{sec:app_assumptions}

Table~\ref{tab:assumptions} lists the twelve modelling assumptions that
most affect the headline KPIs, with the evidence basis for each and the
direction of risk if the assumption proves wrong.
A fuller treatment of each item --- including when assumptions were
revised during the project --- is recorded in \texttt{ASSUMPTIONS.md}
in the project repository.

\begin{table}[htbp]
  \centering
  \caption{Key modelling assumptions, their evidence basis, and the
    direction of error if the assumption is violated.}
  \label{tab:assumptions}
  \small
  \setlength{\tabcolsep}{4pt}
  \begin{tabular}{p{4.8cm} p{4.8cm} p{4.8cm}}
    \toprule
    \textbf{Assumption} & \textbf{Basis} & \textbf{Risk if wrong} \\
    \midrule
    One storage bin per material.  SAP records multiple bins for some
    SKUs; the model assigns each material to a single best-fit slot.
    &
    SAP \texttt{özet} sheet decoded 802 rack bins and 2,872 Kardex
    assignments; for multi-bin materials the nearest bin is used.
    &
    Understates real-warehouse flexibility slightly; the SAP baseline's
    advantage could be marginally lower than reported. \\[4pt]

    Fast-mover manual band: the three lowest rack levels (
    up to $\approx 206$ cm) are reachable without a reach truck.
    &
    Site visit (``first three levels reachable by hand''); corroborated
    by F400 video footage and rack PDF level-height stamps.
    &
    If the real threshold is two levels, manual-band placements are
    over-counted and RT utilization is understated across all policies. \\[4pt]

    F400 timing constants extrapolated to all nine product families.
    Operator pick: 0.113 min; RT pick/place: 0.110 min;
    manual-penalty: 0.102 min; Kardex pick: 0.113 min.
    &
    F400 video time-motion study, 2,319 micro-events, 296.6 min of footage
    (source: \texttt{F400 Kit Cansu Nehir.xlsx}).
    &
    Families with different bin heights or heavier parts may have
    different rates; sensitivity analysis bounds the effect at roughly
    $\pm$15\% on lead time. \\[4pt]

    Service times are lognormal with means from the F400 study and
    log-$\sigma$ moment-matched from the measured CVs
    (operator 1.245, RT 1.047, manual-penalty 1.279, Kardex 1.245).
    &
    F400 CVs: operator pooled CV = 1.93 ($n = 900$); RT CV = 1.41
    ($n = 69$).
    &
    Lognormal may underfit very heavy service tails; if the true
    tail is heavier, RT queue lengths are understated under all
    congestion-prone policies. \\[4pt]

    Batch arrivals: empirical inter-batch gaps and batch sizes sampled
    from the timestamped ZWM92 subset (Okken, AKS\_PAK, F400;
    18,164 orders, 3,885 batches).
    &
    ZWM92 dispatch log; mean batch size 4.68, CV 2.04 (rules out
    Exponential); calibrated at 419 orders/day vs.\ actual 396/day
    ($+5.7\%$, within the $\pm 10\%$ acceptance band).
    &
    Batch sizes are extrapolated from three families to all eight lines;
    if other lines have larger batches, the system is slightly
    under-loaded. \\[4pt]

    Corridor-aware rectilinear routing: travel follows rack-edge safe
    lines; a clear Manhattan L-path is returned immediately; rack
    bodies are treated as obstacles.
    &
    Floor plan from AutoCAD DWG; measured corridor widths: kit 200 cm,
    RT aisles 300--320 cm, perimeter aisles 280--429 cm.
    &
    One-way aisle flow and agent-to-agent congestion are not modelled;
    walk distances are slightly underestimated in high-load scenarios. \\[4pt]

    RT return-to-depot leg: after delivering a pallet the reach truck
    drives back to the depot as a detached process that holds the RT
    resource until arrival.
    &
    DWG floor plan for depot location; standard RT operating procedure
    confirmed at site visit.
    &
    If trucks park near the last delivery point rather than returning to
    depot, RT queue-wait is somewhat overstated; the directional error
    favours Line-aware Slotting (lower RT demand). \\[4pt]

    Kardex single carousel rotation per kit: one rotation is charged
    per kit order regardless of how many Kardex picks that order
    contains; picks inside the rotation are sequential.
    &
    H1 code-audit fix (May 2026); operational semantics of vertical
    carousels confirmed with warehouse staff.
    &
    If each pick triggers an independent rotation, Kardex utilization
    is understated by roughly $3\times$ for kits with multiple Kardex
    items. \\[4pt]

    Single-shift model: ZWM92's four-month dispatch log is compressed
    into an independent 480-minute working day (07:00--16:00 profile).
    Warmup and cooldown windows of 30 min each trim edge effects.
    &
    ZWM92 timestamps span a single shift window; multi-shift dispatch
    patterns were not observed in the log.
    &
    If the real plant runs double shifts, steady-state queue lengths
    could be higher; the model cannot replicate end-of-shift catch-up
    dynamics. \\[4pt]

    Capacity-heuristic placement for undecoded materials (approximately
    2,319 rack materials without a decoded SAP bin): assigned to the
    first free slot in the rack with available capacity.
    &
    750 materials have fully decoded SAP bins; 2,872 are Kardex-routed
    (policy-invariant); the remaining positions use the heuristic.
    &
    The SAP baseline is in practice a SAP+heuristic hybrid; pure SAP
    fidelity is 12.6\% of all active materials (750/5\,941).  Disclosed in the Results
    section. \\[4pt]

    Distribution-driven order arrivals (not trace replay): each
    replication draws fresh samples from fitted distributions with a
    replication-specific seed ($42 + \text{rep}$) and independent
    arrival and service RNG streams.
    &
    Arena-style design per advisor guidance; CRN across policies
    enforced by stream separation.
    &
    Training and validation use the same four-month dataset, so
    the chi-square and volume checks are internal-consistency tests,
    not holdout validation.  Supporting evidence comes from F400 timing,
    daily-volume calibration, and face validity checks in the 3D viewer. \\[4pt]

    Model scope is order picking only; putaway and replenishment are
    out of scope (pallet positions do not deplete).
    &
    Scope agreed with the project supervisor; replenishment scheduling
    was not visible in the ZWM92 dispatch log.
    &
    Line-aware Slotting places fast movers at low (no-RT) levels, which
    would require periodic replenishment in practice.  The policy
    comparison is internally fair --- scope is identical for all six
    policies --- but the RT utilization figures cover picking only. \\
    \bottomrule
  \end{tabular}
\end{table}

% ─────────────────────────────────────────────────────────
\subsection{Software and Reproducibility}\label{sec:app_software}

We wrote the simulation in Python~3 using the SimPy discrete-event
library \cite{simpy}.  The model spans 17 source modules in the
\texttt{src/} directory: data loading and SAP bin decoding,
warehouse geometry and routing, slotting policies, SimPy process logic,
KPI recording, statistical analysis, validation, and sensitivity sweeping.

All statistics in this report are recomputed by script from the
per-replication output files.  Running \texttt{src/main.py} executes the six
policies across 20 replications with common random numbers and writes a
tidy CSV (\texttt{kpi\_by\_replication.csv}) where each row is one
policy--replication pair.  We imported that file directly into Minitab
for the ANOVA and paired $t$-tests in the Results section;
\texttt{src/analyze.py} runs the same tests as a cross-check.
The sensitivity sweeps in \texttt{src/sensitivity.py} apply
one-at-a-time perturbations to each timing constant and produce
the tornado charts in Appendix~\ref{sec:app_sensitivity}.

For face validation we built an interactive 3D warehouse viewer in Three.js.
It renders rack geometry from the AutoCAD DWG and the eleven rack PDF
drawings, animates agent trajectories from the simulation recorder output,
and shows the per-policy KPI comparison panel we used in our project
presentations.

% ─────────────────────────────────────────────────────────
\subsection{Additional Sensitivity Charts}\label{sec:app_sensitivity}

Figure~\ref{fig:sensitivity_2x2} gives OAT tornado charts for three
additional KPIs --- average walking distance, orders completed, and RT
utilization --- alongside the pick-density heatmap for the Usage-based ABC
policy.  The lead-time tornado chart appears in the main Results section.

\begin{figure}[htbp]
  \centering
  \begin{tabular}{cc}
    \includegraphics[width=0.47\textwidth]{sensitivity_tornado_avg_walk_distance} &
    \includegraphics[width=0.47\textwidth]{sensitivity_tornado_orders_completed} \\[6pt]
    \includegraphics[width=0.47\textwidth]{sensitivity_tornado_reach_truck_utilization} &
    \includegraphics[width=0.47\textwidth]{heatmap_usage-based_abc} \\
  \end{tabular}
  \caption{Top row: OAT tornado charts for average walking distance per
    order (left) and orders completed over the five-day run (right).
    Bottom left: OAT tornado chart for reach-truck utilization, showing
    that the arrival rate and RT lift time are the dominant
    drivers of RT queue load.
    Bottom right: pick-density heatmap for the Usage-based ABC policy,
    showing how that policy distributes frequent picks across the rack
    array; the concentration in the central racks contrasts with the
    line-corridor grouping achieved by Line-aware Slotting.}
  \label{fig:sensitivity_2x2}
\end{figure}
